\chapter{Sistemas Distribuidos para \emph{internet of things}}

\section{Familiarización con sockets}
Tome como referencia el archivo TCPServer.c y modifiquelo para que atienda a varios clientes de forma simultáneamente mediante la implementación utilizando la primitiva \emph{select}. 
\section{Protocolos:máquinas de estados}

La empresa QUIJOTECOFFE S.L ha decidido dar un paso mas en su proceso de I+D
y diseñar máquinas expendedoras de café cuya interacción será a través de una App Android. La máquina no tendrá ningún tipo de botón físico/ranura para dinero, etc. sólo y exclusivamente el hueco por donde proporciona el café al usuario. La máquina está dotada de un módem Bluetooth con el que interacciona con el usuario y de un módem GPRS con el que informa al proveedor de mantenimiento acerca del estado de los suministros (café, leche, azúcar, etc.) y estadísticas globales. Como ingeniero de diseño de I+D usted es el encargado del diseño del protocolo de comunicaciones y la máquina de estados.
\begin{itemize}
\item Identifique las entidades del sistema y su relación
\item Especifique los protocolos entre entidades y sus servicios
\item Diseñe la estructura de los mensajes con la herramienta protobuff de google.
\item Dibuje la máquina de estados de la máquina de café. 
\end{itemize}

\section{RMI: Cliente y Servidor Básico en Ice}
\begin{itemize}
\item Ejecutar slice2py iot.ice
\item Implementar en un archivo denominado servidor.py un servidor para la interfaz DigitalSink
\item Implementar un cliente denominado cliente.py que invoque la interfaz update del servidor DigitalSink que usted ha implementado rellenando solo el valor que actualiza.
\item Pruebe su cliente con el servidor implementado por sus compañeros cambiando el endpoint.
\item Pruebe su servidor con los clientes implementados por sus compañeros proporcionándoles el endpoint. 
\item Capture el tráfico y muestre el dialogo TCP de la invocación con WireSark,enumere los métodos del objeto que se invocan. 
\item ¿Qué muestra la herramienta nmap del puerto 10000 donde se ejecuta el servidor?
\item Intente conectarse al object adapter del servidor por medio del netcat, ¿qué contesta el object adapter?.
  \item Haga una invocación completa (añadiendo metadatos).
\end{itemize}

\section{Comunicación indirecta: IceStorm}

Copie su implementación anterior y modifiquela para que se use el
servicio IceStorm.

\section{Diseño de interfaces y gestión}
Utilizando el middleware \textbf{ZeroC Ice} para el ejemplo de la máquina QUIJOTECOFFE S.L.  
\begin{itemize}
	\item Diseñe un interfaz slice para el protocolo definido.
	\item Utilizando las herramientas de la ZeroC Ice y para la parte ejecutándose en la máquina expendedora de café, describa brevemente la forma de configurar el sistema para poder actualizar dicho software de forma remota. Por ejemplo, para 100 máquinas conectadas. 
\end{itemize}


\section{Coordinación}
Utilizando el middleware \textbf{ZeroC Ice} y con el concepto de exclusión mutua basado en servidor central  
\begin{itemize}
	\item Diseñe un interfaz slice para el algoritmo junto con uno/varios compañeros.
	\item Implemente un cliente/servidor para dicho algoritmo. Indique qué compañero implementa el servidor/cliente para su elección. Es decir, si usted implementa un cliente, indique quién implementa el servidor correspondiente.
	\item Pruebe varios clientes/servidores. 
\end{itemize}
